\documentclass[10pt,a4paper]{article}
\usepackage[utf8]{inputenc}
\usepackage[a4paper,left=3cm,right=2cm]{geometry}
\usepackage{amsmath}
\usepackage[thmmarks, amsmath, thref]{ntheorem}
\usepackage{amsfonts}
\usepackage{amssymb}
\usepackage{fancyhdr}
\pagestyle{fancy}
\usepackage{anysize}
\usepackage{graphicx}
\usepackage{float}
\usepackage{hyperref}
\usepackage{enumitem}
\usepackage{verbatim}
\usepackage{siunitx}
\newtheorem{theorem}{Theorem}
\theoremsymbol{\ensuremath{\blacksquare}}
\newtheorem*{solution}{Solución:}

\usepackage{xcolor}

% Margenes y formalidades

\textwidth 6.25in % Margen Derecho y ancho del texto
\textheight 8.5in
\oddsidemargin 0in % Margen Izquierdo
\headheight 0in % Largo del texto
\leftmargin 1in
\parindent 0em % Salto entre la indentación (Sangría)https://www.overleaf.com/4519127669sdpptkhvsbfn
\topmargin -0.5in % Margen Superior
\parskip 0.5ex % Salto entre parrafos
\baselineskip 1pt



\lhead{
\includegraphics[scale=0.08]{Figuras/logoUSM.png}} % texto izquierda de la cabecera
%\chead{}                                      % texto centro de la cabecera
\rhead{Universidad Técnica Federico Santa María\\Ingeniería Comercial } % número de página a la derecha
\renewcommand{\headrulewidth}{0.4pt} % grosor de la línea de la cabecera
\renewcommand{\footrulewidth}{0.4pt} % grosor de la línea del pie

\begin{document}
\vspace{5cm}



\begin{center}
\begin{large}
\textbf{\\Control 3}
\end{large}\\
\end{center}

\begin{center}
\begin{large}
\textbf{Econometría ICS-294}
\end{large}\\
\vspace{0.2cm}
III Trimestre 2025\\
Profesor: Pedro Comte\\
%Ayudante: 
\end{center}

\textbf{Instrucciones:} 

\begin{itemize}

    \item Usted dispone de 60 minutos  para intentar los 100 puntos del control.
    \item Cuide la presentación y redacción de sus preguntas.
    \item No puede utilizar celular, computador o tablet, pero sí una calculadora TI.
     
\end{itemize}

\section{Preguntas cortas (40 puntos)}

\begin{enumerate}



\item (20 puntos) 
Se presenta el siguiente modelo econométrico para explicar las importaciones (IMP) en función del Producto Interno Bruto (PIB), la Inversión (INV) y el Consumo (CONS):

\[
\text{IMP} = -19.7251 + 0.0322 \cdot \text{PIB} + 0.4142 \cdot \text{INV} + 0.2427 \cdot \text{CONS}
\]

Considerando la información disponible, responda lo siguiente:

\begin{enumerate}
    \item (10 puntos) ¿Existe evidencia de multicolinealidad entre las variables explicativas? Fundamente su respuesta.

Para su respuesta considere la siguiente información.

    Matriz de correlaciones:

    \[
    \begin{array}{c|ccc}
    & \text{PIB} & \text{INV} & \text{CONS} \\
    \hline
    \text{PIB} & 1 & 0.2154 & 0.9989 \\
    \text{INV} & 0.2154 & 1 & 0.2137 \\
    \text{CONS} & 0.9989 & 0.2137 & 1 \\
    \end{array}
    \]

    Coeficientes de determinación de regresiones auxiliares para cada regresor:

    \[
    R_{PIB}^2 = 0.9979, \quad R_{INV}^2 = 0.0475, \quad R_{CONS}^2 = 0.9979
    \]



{\color{black}
    \item (10 puntos) En caso de confirmarse la existencia de multicolinealidad, ¿qué métodos o estrategias podrían emplearse para mitigar este problema en el modelo?
\end{enumerate}

{\color{black}
\item (20 puntos)
¿Cómo podría utilizar el test de White (simplificado) para detectar heterocedasticidad en el modelo \( y_i = \beta_0 + \beta_1 x_{1i} + \beta_2 x_{2i} + u_i \)? Describa paso a paso el procedimiento e indique cómo se interpreta el resultado
\end{enumerate}

\color{black}
\section{Ejercicios (60 puntos)}


\begin{enumerate}

\item (30 puntos) A continuación se presenta el gráfico de la serie temporal del tipo de cambio nominal (pesos por dólar) en Chile, correspondiente al período comprendido entre enero de 1983 y diciembre de 2024, con datos de frecuencia mensual. 


\begin{center}
\textbf{Figura 1: Tipos de cambio (pesos por dólar)}
\begin{figure}[H]
    \centering
    \includegraphics[width=0.85\linewidth]{certamenes/BCChChart.png}
\end{figure}
\end{center}

Con el objetivo de evaluar si la serie temporal sigue un proceso de caminata aleatoria (random walk), un grupo de investigadores aplica el test de Dickey-Fuller mediante la estimación de dos modelos auxiliares:

   \textbf{Modelo sin tendencia (without trend):}
    \begin{table}[H]
    \centering
    \begin{tabular}{lccc}
    \toprule
    Variable & Coef. & Error estándar & Estadístico t \\
    \midrule
    Intercepto & 0,0710 & 0,0354 & 2,006 \\
    $y_{t-1}$  & -0,0182 & 0,0097 & -1,876 \\
    \bottomrule
    \end{tabular}
    \end{table}

    \textbf{Modelo auxiliar con tendencia (\textit{trend}):}
    \begin{table}[H]
    \centering
    \begin{tabular}{lccc}
    \toprule
    Variable & Coef. & Error estándar & Estadístico t \\
    \midrule
    Intercepto & 0,0452 & 0,0289 & 1,563 \\
    Tendencia ($t$) & 0,0021 & 0,0003 & 7,000 \\
    $y_{t-1}$  & -0,0304 & 0,018 & -1,6889 \\
    \bottomrule
    \end{tabular}
    \end{table}

  Testee con un 1\% de significancia si el tipo cambio nominal tiene raíz unitaria. Explique su resultado.  




\item (30 puntos) Suponga que usted está interesado en determinar los factores que explican
las diferencias salariales entre distintos individuos. Para realizar
su análisis, usted especifica la siguiente función de regresión:

\begin{equation}
wage = \beta_{1} + \beta_{2}IQ + \beta_{3}Educ + \beta_{4}Exper + \beta_{5}Tenure + \beta_{6}personacolor + u \label{eq:original}
\end{equation}

donde \emph{wage} representa el salario, \emph{IQ} es el coeficiente intelectual,
\emph{Educ} son los años de educación, \emph{Exper} los años de experiencia laboral,
\emph{Tenure} la antigüedad en el trabajo actual, y \emph{personacolor} es una variable dicotómica
que toma el valor 1 cuando la persona se identifica como {persona de color}, y 0 en caso contrario.


Los resultados de la regresión se presentan en la Tabla 1, y en la Figura 1 se muestran los residuos
agrupados según la variable \emph{personacolor}:

\begin{center}
\textbf{Tabla 1}
\end{center}

\begin{center}
Model 1: OLS, using observations 1--935\\ Dependent variable: wage\\
\vspace{1em}
\begin{tabular}{lr@{.}lr@{.}lr@{.}lr@{.}l}
 &  \multicolumn{2}{c}{Coefficient} &   \multicolumn{2}{c}{Std.\ Error} &    \multicolumn{2}{c}{$t$-ratio} &     \multicolumn{2}{c}{p-value} \\
\hline
const &   $-$405&803 &     122&400 &       $-$3&3154 &         0&0010 \\
IQ &   3&77957 &     0&996294 &       3&7936 &         0&0002 \\
educ &   57&5438 &     7&00006 &       8&2205 &         0&0000 \\
exper &   14&8737 &     3&19330 &       4&6578 &         0&0000 \\
tenure &   6&88539 &     2&46227 &       2&7964 &         0&0053 \\
personacolor & -40&15 & 38&70 & -1&037 & 0&296\\
\end{tabular}
\vspace{1ex}

\begin{tabular}{lrlr}
Mean dependent var &  957.9455 & S.D. dependent var &  404.3608 \\
Sum squared resid &  1.25\textrm{e+08} & S.E. of regression &  367.4213 \\
$R^2$ &  0.178780 & Adjusted $R^2$ &  0.174361 \\
$F(5, 929)$ &  40.44887 & P-value($F$) &  1.07\textrm{e--37} \\
\end{tabular}
\end{center}

\begin{center}
\textbf{Figura 1: Residuos agrupados por identificación étnico-racial}
\begin{figure}[H]
    \centering
    \includegraphics[width=0.65\linewidth]{certamenes/racial.png}
\end{figure}
\end{center}

\begin{enumerate}
    \item (6 puntos) De acuerdo a esta información, ¿le parece que existe heterocedasticidad en el modelo? ¿Tiene sentido este resultado? Explique.

 

    \item (10 puntos)  En la Tabla 2 se presentan los resultados de una regresión de los residuos al cuadrado de la regresión de la tabla 1 sobre las variables explicativas:

\begin{center}
\textbf{Tabla 2: OLS, usando 935 observaciones} \\
Variable dependiente: $\hat{u}^2$ \\
\vspace{1em}
\begin{tabular}{lr@{.}lr@{.}lr@{.}lr@{.}l}
 & \multicolumn{2}{c}{Coeficiente} & \multicolumn{2}{c}{Error estándar} & \multicolumn{2}{c}{$t$-ratio} & \multicolumn{2}{c}{p-valor} \\[1ex]
const          & $-$160277&000 & 104551&00 & $-$1&5330 & 0&1256 \\
IQ             & 921&598   & 851&010  & 1&0829  & 0&2791 \\
educ           & 16288&600 & 5979&280 & 2&7242  & 0&0066 \\
exper          & 2955&490  & 2727&640 & 1&0835  & 0&2789 \\
tenure         & $-$6586&500 & 2103&220 & $-$3&1316 & 0&0018 \\
personacolor   & $-$70000&000 & 30000&000 & $-$2&3333 & 0&0200 \\
\end{tabular}

\vspace{1ex}
\begin{tabular}{lrlr}
Media variable dependiente & 134132.1 & Desv. estándar dependiente & 317721.3 \\
Suma de residuos al cuadrado & $9.15 \times 10^{13}$ & Error estándar de regresión & 313842.4 \\
$R^2$ & 0.029491 & $R^2$ ajustado & 0.024268 \\
$F(5, 929)$ & 5.646001 & Valor-p ($F$) & 0.000039 \\
\end{tabular}
\end{center}


    \[
    \hat{\mu}^2 = a_1 + a_2 IQ + a_3 Educ + a_4 Exper + a_5 Tenure + a_6 personacolor + v
    \]

    ¿Puede usarse esta regresión para testear la presencia de heterocedasticidad? Explique e interprete.

      \item (8 puntos) Observe que la variable \texttt{personacolor} no es estadísticamente significativa según los resultados de la tabla 1. ¿Podemos concluir que \texttt{personacolor} no afecta el ingreso? Fundamente su respuesta considerando la presencia de heterocedasticidad. 

        \item (6 puntos) ¿Cómo se puede corregir el problema anterior para obtener inferencias válidas? ¿Qué tipo de errores estándar se deben utilizar? 

   
\end{enumerate}

\end{enumerate}














\newpage


\begin{center}
    {\Large Formulario  }

 
\end{center}

\begin{itemize}

% \item  
% {\color{red}
% $$\hat{\beta}_0 = \bar{y} - \hat{\beta}_1 \bar{x} $$
% $$\hat{\beta}_1 = \frac{\sum_{i=1}^{n}(x_i - \bar{x})(y_i-\bar{y})}{\sum_{i=1}^{n}(x_i - \bar{x})^2}$$}

\item  {\color{red} \[ \hat{\beta} = (X^{\top}X)^{-1}X^{\top}y \]}

\item 
	{\color{red}	\[ SST = \sum_{i=1}^{n} (y_i - \bar{y})^2 \quad SSE = \sum_{i=1}^{n} (\hat{y}_i - \bar{y})^2  \quad SSR = \sum_{i=1}^{n} \hat{u}_i^2\]}
		

    \item {\color{red} \[ R^2  = 1 - \frac{SSR}{SST} \quad \quad \bar{R}^2 =  1 - \frac{(n-1)SSR}{(n-k-1)SST} \]}

 \item {\color{red} \[ t_{\hat{\beta_j}}=\frac{\hat{\beta}_j - \beta_j}{\text{s.e.}(\hat{\beta}_j)} \sim t_{n-k-1} \]}

 \item {\color{red}\[
    F = \frac{SSR_{r} - SSR_{ur}}{SSR_{ur}} \cdot \frac{n - k - 1}{j} \sim F(j, n - k - 1)
    \]}

\item {\color{red}Test F para regresión auxiliar.}
\begin{center}
 {  \color{red}     \(F=\dfrac{R^{2}/j}{(1-R^{2})/(n-k-1)}
          \sim F(j,\;n-k-1)\).}
\end{center}
% \item {\color{red}Considere que el valor crítico a un nivel de significancia del 5\% es 1.645}
% \item {\color{red}Considere los siguientes valores críticos para t: $\alpha = 0.01 \implies t \approx 2.62$; $\alpha = 0.05 \implies t \approx 1,97$;$\alpha = 0.01 \implies t \approx 2.62$; $\alpha = 0.1 \implies t \approx 1,645$}
%\item {\color{red}Considere que el valor crítico de F al 1\% con grados de libertad (2,178) es aproximadamente 4.72}


\item \color{red}{ Multiplicador de Lagrange

$$LM = nR^2 \sim \chi^2_{df}$$

Tabla $\chi^2_{df}$
\begin{center}
    \begin{tabular}{|c|c|c|c|}
    \hline
        df & 0.99 & 0.95 & 0.90 \\ 
        \hline
        1 &  6.635 & 3.841 & 2.706\\
        \hline
        2 &  9.210 & 5.991 & 4.605 \\
        \hline
        3 & 11.345 & 7.815 & 6.251  \\
        \hline
        4 &  13.277 & 9.488 & 7.779 \\
        \hline
        5 &  15.086 & 11.070 & 9.236 \\
        \hline
        6 &  16.812 & 12.592 & 10645 \\
        \hline
    \end{tabular}
    \end{center}}


\item Tabla de test D-F
\begin{center}
    \begin{figure}[H]
        \centering
        \includegraphics[width=0.7\linewidth]{sem14/tabla.png}
    \end{figure}
\end{center}

\end{itemize}


% \begin{figure}[H]
%     \centering
%     \includegraphics[width=0.85\linewidth]{Ejercicios//ejercicios_inferencia/cert2_F2_1184.png}
% \end{figure}

% \begin{figure}[H]
%     \centering
%     \includegraphics[width=0.85\linewidth]{certamenes/normal.png}
% \end{figure}


\begin{figure}
    \centering
    \includegraphics[width=0.85\linewidth]{certamenes/normal_1cola.png}
\end{figure}



\end{document}